%% toc: P2 19 19.8
\section{Proofs for factoring under translation and superposition into either side}
\label{sec:annex-p02-w02c-p02-021}

The translation, the fresh symbol \(\mathsf{eq}\) and the facts (F0) and (F1)
are those of \cref{sec:annex-p02-w02b-p02-026}.  Two consequences of
Lemma~\ref{thm:p02-w02b-p02-026-translation}(1) are used below.

\emph{(T1) Polarity and atoms.}  For an \(\Lang\)-atom \(A\),
\(A^{\mathsf{eq}}\) is an atom and \((\neg A)^{\mathsf{eq}}=\neg(A^{\mathsf{eq}})\).
Hence a literal \(L\) and its translation have the same polarity, and the
atom of \(L^{\mathsf{eq}}\) is the translation of the atom of \(L\).

\emph{(T2) Injectivity on clauses.}  Since translation is injective on
literals, \(L\in C\) iff \(L^{\mathsf{eq}}\in C^{\mathsf{eq}}\) for every
literal \(L\) and clause \(C\).  So if \(C^{\mathsf{eq}}=D^{\mathsf{eq}}\),
then for every literal \(L\), \(L\in C\) iff \(L^{\mathsf{eq}}\in
C^{\mathsf{eq}}=D^{\mathsf{eq}}\) iff \(L\in D\); that is, \(C=D\).  Likewise
\(L_1\neq L_2\) iff \(L_1^{\mathsf{eq}}\neq L_2^{\mathsf{eq}}\).

\begin{annexitem}{annex:p02-w02c-p02-021-factoring}{thm:p02-w02c-p02-021-factoring}{Factoring commutes with translation}
\begin{proof}
(1)  Let \(L_1,L_2\in C\) be distinct literals of the same polarity with atoms
\(A\) and \(B\), let \(\theta\) be a most general unifier of \(A\) and \(B\),
and let \(C'=C\theta\).  Then \(L_1^{\mathsf{eq}},L_2^{\mathsf{eq}}\) belong to
\(C^{\mathsf{eq}}\) by the definition of \(C^{\mathsf{eq}}\); they are
distinct by (T2); they have the common polarity of \(L_1\) and \(L_2\), and
their atoms are \(A^{\mathsf{eq}}\) and \(B^{\mathsf{eq}}\), by (T1).  By
Lemma~\ref{thm:p02-w02b-p02-026-translation}(2), \(\theta\) is a most general
unifier of \(A^{\mathsf{eq}}\) and \(B^{\mathsf{eq}}\).  So Factoring applies
to \(C^{\mathsf{eq}}\) on \(L_1^{\mathsf{eq}},L_2^{\mathsf{eq}}\) with
\(\theta\), and its conclusion is \(C^{\mathsf{eq}}\theta\), which equals
\((C\theta)^{\mathsf{eq}}=C'^{\mathsf{eq}}\) by
Lemma~\ref{thm:p02-w02b-p02-026-translation}(1).

(2)  Let \(M_1,M_2\in C^{\mathsf{eq}}\) be distinct literals of the same
polarity with atoms \(A'\) and \(B'\), let \(\theta\) be a most general
unifier of \(A'\) and \(B'\), and let \(C'^{\mathsf{eq}}=C^{\mathsf{eq}}\theta\).

\emph{The literals.}  As \(C^{\mathsf{eq}}=\{L^{\mathsf{eq}}:L\in C\}\), there
are \(L_1,L_2\in C\) with \(L_1^{\mathsf{eq}}=M_1\) and
\(L_2^{\mathsf{eq}}=M_2\); by injectivity on literals they are unique.  They
are distinct, since \(L_1=L_2\) would give \(M_1=M_2\).  By (T1), \(L_i\) has
the polarity of \(M_i\) for \(i=1,2\), so \(L_1\) and \(L_2\) have the same
polarity.

\emph{The unifier.}  Let \(A\) and \(B\) be the atoms of \(L_1\) and \(L_2\).
By (T1), the atoms of \(M_1\) and \(M_2\) are \(A^{\mathsf{eq}}\) and
\(B^{\mathsf{eq}}\), so \(A'=A^{\mathsf{eq}}\) and \(B'=B^{\mathsf{eq}}\).  By
Lemma~\ref{thm:p02-w02b-p02-026-translation}(2), \(\theta\) is a most general
unifier of \(A'\) and \(B'\) iff it is one of \(A\) and \(B\); so it is a most
general unifier of \(A\) and \(B\).

\emph{The conclusion.}  Factoring therefore applies to \(C\) on \(L_1,L_2\)
with \(\theta\), and its conclusion is \(C\theta\).  By
Lemma~\ref{thm:p02-w02b-p02-026-translation}(1),
\((C\theta)^{\mathsf{eq}}=C^{\mathsf{eq}}\theta=C'^{\mathsf{eq}}\), and by
(T2), \(C\theta=C'\).  So \(C'\) follows from \(C\) by Factoring on
\(L_1,L_2\) with \(\theta\).

\emph{Item (3) of the translation lemma.}  Let \(T\) be a set of
\(\Lang\)-clauses.  Let \(C_1,\ldots,C_n\) be \(\Lang\)-clauses such that
\(C_1^{\mathsf{eq}},\ldots,C_n^{\mathsf{eq}}\) is a first-order resolution
derivation from \(T^{\mathsf{eq}}\), and let \(k\leq n\).  There are three
cases.  If \(C_k^{\mathsf{eq}}\in T^{\mathsf{eq}}\), then
\(C_k^{\mathsf{eq}}=D^{\mathsf{eq}}\) for some \(D\in T\), and \(C_k=D\in T\)
by (T2).  If \(C_k^{\mathsf{eq}}\) follows by Resolution from renamed-apart
copies of earlier members \(C_i^{\mathsf{eq}}\) and \(C_j^{\mathsf{eq}}\),
then \(C_k\) follows by Resolution from renamed-apart copies of \(C_i\) and
\(C_j\), by the converse Resolution argument in the proof of
Lemma~\ref{thm:p02-w02b-p02-026-translation}.  If \(C_k^{\mathsf{eq}}\)
follows by Factoring from an earlier member \(C_i^{\mathsf{eq}}\), on two
literals with a most general unifier \(\theta\), then \(C_k\) follows from
\(C_i\) by Factoring with \(\theta\), by item~(2) above.  So
\(C_1,\ldots,C_n\) is a first-order resolution derivation from \(T\).  The
direction from \(\Lang\)-derivations to translated derivations is the
forward argument in the proof of
Lemma~\ref{thm:p02-w02b-p02-026-translation} for members of \(T\) and for
Resolution, together with item~(1) above for Factoring.  This proves
item~(3) in both directions, and with it the steps of the proofs of
Theorems~\ref{thm:p02-w02b-p02-026-eq-resolution}
and~\ref{thm:p02-w02b-p02-026-bare} that read every derivation from a
translated set as the translation of a derivation from the untranslated set.
\end{proof}
\end{annexitem}

\begin{annexitem}{annex:p02-w02c-p02-021-superposition-rhs}{thm:p02-w02c-p02-021-superposition-rhs}{Superposition into the second argument is sound}
\begin{proof}
Positions, \(s|_p\), \(s[u]_p\) and the facts (P1) and (P2) are those
established at the start of the proof of
Lemma~\ref{thm:p02-w02b-p02-026-superposition-sound}: (P1) if \(p\) is a
position of \(s\), then \(p\) is a position of \(s\sigma\),
\((s\sigma)|_p=(s|_p)\sigma\) and \((s[u]_p)\sigma=(s\sigma)[u\sigma]_p\); (P2)
if \(u\) and \(u'\) have the same value in \(\mathcal N\) under an assignment
\(\beta\), then so do \(v[u]_p\) and \(v[u']_p\) for every position \(p\) of
\(v\).

\emph{Setup.}  Let \(\mathcal N\) satisfy \(C\cup\{K\}\) and
\(D\cup\{t\mathrel{\diamond}s\}\).  By Lemma~\ref{lem:clause-substitution},
applied to each premise separately, \(\mathcal N\) satisfies
\((C\cup\{K\})\sigma\) and \((D\cup\{t\mathrel{\diamond}s\})\sigma\).  Fix an
assignment \(\beta\).  By (F0) it suffices to find a literal of the conclusion
true in \(\mathcal N\) under \(\beta\).  Every literal of \(C\sigma\) or
\(D\sigma\) belongs to the conclusion, so if \(\beta\) satisfies one of them
we are done.  Assume therefore that it satisfies none.  By (F0) applied to the
two instantiated premises, \(\beta\) satisfies \(K\sigma\) and
\((t\mathrel{\diamond}s)\sigma\).

\emph{The equation.}  \(K\sigma\) is \(l\sigma=r\sigma\) or
\(r\sigma=l\sigma\).  In either case \(l\sigma\) and \(r\sigma\) have the
same value under \(\beta\), since \(=\) is interpreted as identity.

\emph{The replaced term (P1).}  Since \(s|_p=l'\), the term \(s\) is
\(s[l']_p\).  By (P1), \(p\) is a position of \(s\sigma\), and
\[
 s\sigma=(s[l']_p)\sigma=(s\sigma)[l'\sigma]_p=(s\sigma)[l\sigma]_p,
 \qquad
 (s[r]_p)\sigma=(s\sigma)[r\sigma]_p,
\]
where the third equality uses \(l'\sigma=l\sigma\).

\emph{Equal values (P2).}  Apply (P2) with \(v=s\sigma\), the position \(p\)
of \(s\sigma\), \(u=l\sigma\) and \(u'=r\sigma\), which have the same value
under \(\beta\).  It gives that \((s\sigma)[l\sigma]_p\) and
\((s\sigma)[r\sigma]_p\), that is \(s\sigma\) and \((s[r]_p)\sigma\), have the
same value under \(\beta\).

\emph{The target literal.}  Substitution acts on both arguments of an atom
and commutes with negation, so \((t\mathrel{\diamond}s)\sigma\) is
\(t\sigma\mathrel{\diamond}s\sigma\) and
\((t\mathrel{\diamond}s[r]_p)\sigma\) is
\(t\sigma\mathrel{\diamond}(s[r]_p)\sigma\).  The truth of a literal
\(w\mathrel{\diamond}w'\) under \(\beta\) depends only on the values of \(w\)
and \(w'\): when \(\diamond\) is \(=\) it is true iff the values coincide, and
when \(\diamond\) is the negation of \(=\) it is true iff they differ.  The
literals \(t\sigma\mathrel{\diamond}s\sigma\) and
\(t\sigma\mathrel{\diamond}(s[r]_p)\sigma\) have the same first argument and
second arguments of equal value, so the second is true under \(\beta\)
because the first is.  It belongs to the conclusion.

\emph{The first argument.}  If the target literal is \(s\mathrel{\diamond}t\)
and the conclusion contains \(s[r]_p\mathrel{\diamond}t\), the steps (P1) and
(P2) are unchanged, and the last step compares
\(s\sigma\mathrel{\diamond}t\sigma\) with
\((s[r]_p)\sigma\mathrel{\diamond}t\sigma\), which have first arguments of
equal value and the same second argument; so again the conclusion literal is
true under \(\beta\).

In every case \(\beta\) satisfies a literal of the conclusion.  As \(\beta\)
was arbitrary, \(\mathcal N\) satisfies the conclusion by (F0).  Only
\(l\sigma=l'\sigma\) was used; no ordering, eligibility or selection
condition entered.
\end{proof}
\end{annexitem}
